% !TEX program = pdflatex

%%%%%%%%%%%%%%%%%%%%
% affineConwayCompData.tex
%%%%%%%%%%%%%%%%%%%%%

\pdfoutput=1 %When you use pdflatex
\documentclass{amsart} %When you use pdflatex
%\documentclass[dvipdfmx]{amsart} %When you use uplatex + dvipdfmx


\usepackage{tikz}

\usetikzlibrary{calc}
\usetikzlibrary{decorations.markings}
\usetikzlibrary {arrows.meta}
\usetikzlibrary {bending}
\usetikzlibrary{intersections}

\usepackage{amssymb}
\usepackage{amsmath}
\usepackage{url}
\usepackage{enumerate}
\usepackage{graphicx}

\usepackage{hyperref}
\hypersetup{
 colorlinks = true,
 linkcolor = blue,
 filecolor = blue,
 urlcolor = blue,
 citecolor = blue,
}
\hypersetup{unicode=true}

%\newcommand{\secK3}{\texorpdfstring{$K3$}{K3}}
\renewcommand{\baselinestretch}{1.4}

\usepackage{bm}
%\usepackage{enumerate}
\usepackage{enumitem}

\input mymacros
\input mymacrosaffineConway

\newcommand{\ttbox}[1]{\makebox[0pt][l]{{\tt #1}}} 
\newcommand{\ttfbox}[1]{\makebox[0pt][l]{\fbox{\tt #1}}} 
\newcommand{\pt}{p{.2cm}}

%%%%%%%%%%%%%%%%%%%%%%%%%%%%

\begin{document}

\title[[Involutions in the affine Conway group: Computational data]
{Involutions in the affine Conway group: Computational data}


\author{Ichiro Shimada}
\address{Department of Mathematics,
Graduate School of Science,
Hiroshima University,
1-3-1 Kagamiyama,
Higashi-Hiroshima,
739-8526 JAPAN}
\email{ichiro-shimada@hiroshima-u.ac.jp}
\thanks{Supported by JSPS KAKENHI Grant Number~20H00112, 23K20209, and 23H00081}

%超平面配置に関連する離散構造の拡張、深化とその応用
%23H00081
%代数幾何学の計算機による研究の新展開
%23K20209
%格子、保型形式とK3曲面、エンリケス曲面の研究
%20H00112
\maketitle
%
%
% 
In the  files
\begin{itemize}
\item[]  \tname{affineConwayCompdata.txt} (about 700KB), 
\end{itemize}
 we present  the computational data used in the preprint 
\begin{itemize}
\item []
Ichiro Shimada.  Involutions in the affine {C}onway group. Preprint.
  \end{itemize}
This document explains the structure and contents of the data.
%
\par
{\bf Remarks}
\begin{itemize}
\item
In the following, we freely use the notation introduced in the preprint above.
\item
We use the format of {\tt GAP}.
In particular, 
some parts of the data are given in  the {\tt Record} format of {\tt GAP}.
\item
Elements of a vector space are given by \emph{row} vectors.
A linear map is expressed by the multiplication of a matrix from the \emph{right}.
\end{itemize}
%
\section{Contents of \tname{affineConwayCompdata.txt}}
%
In the following, we arrange the elements of $\Omega=\PP^1(\FF_{23})$ in order as 
\[
\infty, 0, 1,2 , \dots, 22, 
\]
and we  fix isomorphisms  
$\FF_2^{\Omega}\cong \FF_2^{24}$ and  $\ZZ^{\Omega}\cong \ZZ^{24}$
accordingly.
For example, the vector $v_{\infty}\in \ZZ^{\Omega}$ is $[1, 0, \dots, 0]$
and $v_{22}\in \ZZ^{\Omega}$ is $[0, \dots, 0, 1]$.
%
\begin{itemize}
%
\item \tname{theOmega} is the list
$\tt{[infinity, 0,1,2,\dots, 22]}$.
%
 %
 \item \tname{GolayWords} is the extended binary Golay code $\Golay\subset \FF_2^{24}$.
%constructed by the method in~\cite[Chapter~10.2]{TheCSbook}.
This data is the list of $2^{12}=4096$ codewords of $\Golay$,
each of which 
is expressed as a subset of \tname{theOmega}.
%
\item \tname{LeechBasis} gives a basis of the Leech lattice $\Lambda\subset \ZZ^{\Omega}$.
%constructed by the method in~\cite[Chapter~10.3]{TheCSbook}.
This data 
is a list of $24$ vectors in $\ZZ^{\Omega}$.
%
\item \tname{GramLeech} is the Gram matrix of  $\Lambda$
with respect to \tname{LeechBasis}.
%
\item \tname{erecs} is the list of $4$ records,
each of which describes the involution $\vep_n=\vep(C_n)\in \dotz$
for $n=8, 12, 16, 24$.
Each record \tname{erec} has the following structure:
\par
\[
{
\renewcommand{\arraystretch}{.7}
\begin{tabular}{\pt \pt \pt \pt}
\ttfbox{erecs} &&& \\
& \ttbox{n} && \\
& \ttbox{word} && \\
& \ttbox{invol} && \\
%
& \ttbox{plusrec} && \\
&& \ttbox{basis} & \\
&& \ttbox{Gram} & \\
&& \ttbox{OGrec} & \\
&&& \ttbox{size}  \\
&&& \ttbox{gens}  \\
&& \ttbox{Kerqrec} & \\
&&& \ttbox{size}  \\
&&& \ttbox{gens}  \\
%
& \ttbox{minusrec} && \\
&& \ttbox{basis} & \\
&& \ttbox{Gram} & \\
&& \ttbox{OGrec} & \\
&&& \ttbox{size}  \\
&&& \ttbox{gens}  \\
&& \ttbox{Kerqrec} & \\
&&& \ttbox{size}  \\
&&& \ttbox{gens}  \\
\end{tabular}
}
\]
%
%
\begin{itemize}
%
\item \tname{n} is $n\in \{8, 12, 16, 24\}$.
%
\item \tname{word} is the  element $C_n$ of \tname{GolayWords} such that $\vep_n=\vep(C_n)$.
%
\item \tname{invol} is the matrix representation of the involusion $\vep_n\in \dotz$
with respect to the basis  \tname{LeechBasis} of $\Leech$.
%
\item \tname{plusrec}  
is a record that describes the invariant part $\Leech(\vep_n, +)$ of $\vep_n$.
When $n=24$, this is a null record, because $\Leech(\vep_{24}, +)$ is of rank $0$.
It contains the following items:
%
\begin{itemize}
\item \tname{basis}
is a basis of $\Leech(\vep_n, +)$,
which is a list of  vectors in $\Leech$.
\item \tname{Gram}
is the Gram matrix of $\Leech(\vep_n, +)$
with respect to the \tname{basis} above.
%
\item \tname{OGrec}
is a record that describes the group finite $\OG(\Leech(\vep_n, +))$.
It contains the following items:
%
\begin{itemize}
\item \tname{size} is the size of $\OG(\Leech(\vep_n, +))$.
\item \tname{gens} is a finite set of generators of $\OG(\Leech(\vep_n, +))$,
which is a set of matrix representations of isometries of $\Leech(\vep_n, +)$ 
with espect to the \tname{basis} above.
\end{itemize}
%
\item \tname{Kerqrec}
is a record that describes the kernel of 
the natural homomorphism  $\etaplus\colon \OG(\Leech(\vep_n, +))\to \OG(q_+)$
to the automorphism group of the discriminant form  $q_+$ of $\Leech(\vep_n, +)$.
It contains the following items:
%
\begin{itemize}
\item \tname{size} is the size of $\Ker\etaplus$.
\item \tname{gens} is a finite set of generators of $\Ker\etaplus$,
which is a set of matrix representations of isometries of $\Leech(\vep_n, +)$ 
with espect to the \tname{basis} above.
\end{itemize}
%
\end{itemize}
%
\item \tname{minusrec}  
is a record that describes the anti-invariant part $\Leech(\vep_n, -)$ of $\vep_n$.
The items in this record are defined in  the same way as those of \tname{plusrec}
with all  ``$+$" being replaced by ``$-$".
When $n=24$,
the item \tname{minusrec.Kerqrec} is  a null record, %{\tt "not available"},
because the discriminant group of $\Leech(\vep_{24}, -)=\Leech$ is trivial.
%
 \end{itemize}
%
\item \tname{GramL26} is the Gram matrix of $L_{26}=U\oplus \Leech(-1)$
with respect to the basis $\weyl_0$,  $\weyl_1$ of $U$ and 
\tname{LeechBasis} of $\Leech(-1)$.
In the following,
we fix this basis of 
$L_{26}=U\oplus \Leech(-1)$
and refer it as \tname{L26Basis}.

%
\item \tname{gammarecs} is the list of $9$ records,
each of which describes the involution $\gamma=(\vep_n, \tau)\in \dotinf$.
Each record \tname{gammarec} has the following structure:
\par
\[
{
\renewcommand{\arraystretch}{.7}
\begin{tabular}{\pt \pt \pt \pt}
\ttfbox{gammarecs} &&& \\
& \ttbox{n} && \\
& \ttbox{epsilon} && \\
& \ttbox{tau} && \\
& \ttbox{invol} && \\
%
& \ttbox{plusrec} && \\
&& \ttbox{basis} & \\
&& \ttbox{Gram} & \\
&& \ttbox{Xitau} & \\
%
& \ttbox{minusrec} && \\
&& \ttbox{basis} & \\
&& \ttbox{Gram} & \\
&& \ttbox{OGrec} & \\
&&& \ttbox{size}  \\
&&& \ttbox{gens}  \\
&& \ttbox{Kerqrec} & \\
&&& \ttbox{size}  \\
&&& \ttbox{gens}  \\
\end{tabular}
}
\]
%
%
\begin{itemize}
%
\item \tname{n} is the $n$ in $\gamma=(\vep_n, \tau)$,
which is  equal to the rank of $\Lts (\gamma, -)$.
%
\item
\tname{epsilon} is the matrix representation of $\vep_n\in \dotz$
with respect to \tname{LeechBasis}.
%
\item 
\tname{tau} is the vector representation of $\tau\in \Leech$
with respect to \tname{LeechBasis}.
%
\item \tname{invol} is the matrix representation of  $\gamma\in \dotinf$
with respect to the basis  \tname{L26Basis} of $\Lts$.
%
\item \tname{plusrec}  
is a record that describes the invariant part $\Lts(\gamma, +)$ of $\gamma$.
It contains the following items:
%
\begin{itemize}
\item \tname{basis}
is a basis of $\Lts(\gamma, +)$,
which is a list of  vectors.
Each element of \tname{basis} is 
written as a vector  of length $26$ with respect to  \tname{L26Basis}.
\item \tname{Gram}
is the Gram matrix of $\Lts(\gamma, +)$
with respect to the \tname{basis} above.
\end{itemize}
%
%When $n\ne 24$, this item \tname{plusrec}  
%also contains the following item:
%
\begin{itemize}
\item \tname{Xitau}
is the set $\Xi_{\tau}$
when $n\ne 24$,
and is {\tt "not available"} when $n=24$.
For $n\ne 24$,
each element of \tname{Xitau} is written as a vector 
of length $24$ with respect to \tname{LeechBasis}.
(The  set $\Xi_{\tau}$ for $n=24$ is not used in the preprint.)
\end{itemize}
%
\item \tname{minusrec}  
is a record that describes the anti-invariant part $\Lts(\gamma, -)$ of $\gamma$.
It  contains the same items as \tname{erec.plusrec} and  \tname{erec.minusrec},
except that,
for the case  $n=24$,
the records {\tt minusrec.OGrec} and {\tt minusrec.Kerqrec}
are null records.
(The data for $n=24$ are not used in the preprint.)
Remark that,
whereas \tname{erec.minusrec} describes a positive-definite lattice,
the record \tname{gammarec.minusrec} describes a  a \emph{negative}-definite lattice $\Lts(\gamma, -)$.

%
 \end{itemize}
%
%
 \end{itemize}
%
\bibliographystyle{plain}
\bibliography{myrefsaffineConway}
%
\end{document}


